\documentclass[aspectratio=169,11pt]{beamer}

\usetheme{default}
\usepackage[T1]{fontenc}
\usepackage{lmodern}
\usepackage{amsmath,amssymb,bm}
\usepackage{booktabs}
\usepackage{array}
\usepackage{xcolor}

% Minimal one-accent palette for a simple academic presentation
\definecolor{navy}{HTML}{1F3A5F}
\definecolor{blue}{HTML}{1F3A5F}
\definecolor{teal}{HTML}{1F3A5F}
\definecolor{orange}{HTML}{1F3A5F}
\definecolor{cream}{HTML}{FFFFFF}
\definecolor{ink}{HTML}{111111}
\definecolor{muted}{HTML}{555555}

\setbeamercolor{normal text}{fg=ink,bg=cream}
\setbeamercolor{frametitle}{fg=navy,bg=cream}
\setbeamercolor{structure}{fg=blue}
\setbeamercolor{alerted text}{fg=orange}
\setbeamercolor{block title}{fg=navy,bg=white}
\setbeamercolor{block body}{fg=ink,bg=white}
\setbeamercolor{block title example}{fg=navy,bg=white}
\setbeamercolor{block body example}{fg=ink,bg=white}
\setbeamercolor{block title alerted}{fg=navy,bg=white}
\setbeamercolor{block body alerted}{fg=ink,bg=white}
\setbeamercolor{footline}{fg=muted,bg=cream}

\setbeamertemplate{navigation symbols}{}
\setbeamertemplate{itemize item}{\color{teal}$\blacktriangleright$}
\setbeamertemplate{itemize subitem}{\color{orange}$\bullet$}
\setbeamertemplate{blocks}[default]
\setbeamertemplate{frametitle}{%
  \vspace{0.35cm}\hspace{0.15cm}\insertframetitle\par
  \vspace{-0.15cm}\hspace{0.15cm}{\color{navy}\rule{1.1cm}{1.2pt}}
}


\newcommand{\keyeq}[1]{%
  \begin{center}
  \setlength{\fboxsep}{9pt}%
  {\color{navy}$\displaystyle #1$}
  \end{center}}

\title{Vector Spaces, Hyperplanes\and the Perceptron}
\subtitle{Module 2 - Mathematics for Computing}
\author{Austin Sabu}
\date{September 2026}

\begin{document}

% 1
\begin{frame}[plain]
  \vspace{0.8cm}
  {\color{navy}\rule{2.0cm}{1.5pt}}\\[0.45cm]
  {\fontsize{30}{34}\selectfont\bfseries\color{navy}Vector Spaces, Hyperplanes\\and the Perceptron}\\[0.35cm]
  {\large\color{blue}Module 2 - Mathematics for Computing}\\[0.8cm]
  {\normalsize\bfseries Austin Sabu}\\
  {\small MTECH in CSE(AIDS)}\\
\end{frame}

% 2
\begin{frame}{The learning path leads to one classifier}
\begin{columns}[T,onlytextwidth]
\column{0.47\textwidth}
\begin{block}{Mathematical foundation}
\begin{itemize}
  \item Algebraic structures
  \item Vectors and vector spaces
  \item Inner products and norms
  \item Angles and orthogonality
\end{itemize}
\end{block}
\column{0.47\textwidth}
\begin{exampleblock}{Machine-learning application}
\begin{itemize}
  \item Lines, planes and hyperplanes
  \item Weight vectors and bias
  \item Perceptron decision rule
  \item Linear classification
\end{itemize}
\end{exampleblock}
\end{columns}
\vspace{0.25cm}
\keyeq{\hat y=\phi(\bm w^{\mathsf T}\bm x+b)}
\end{frame}

% 3
\begin{frame}{Each algebraic structure adds one property}
\small
\begin{tabular}{>{\bfseries}lccccc}
\toprule
Structure & Closure & Associative & Identity & Inverse & Commutative\\
\midrule
Magma         & \checkmark &--&--&--&--\\
Semigroup     & \checkmark &\checkmark&--&--&--\\
Monoid        & \checkmark &\checkmark&\checkmark&--&--\\
Group         & \checkmark &\checkmark&\checkmark&\checkmark&--\\
Abelian group & \checkmark &\checkmark&\checkmark&\checkmark&\checkmark\\
\bottomrule
\end{tabular}
\vspace{0.5cm}
\begin{itemize}
  \item Closure: $a,b\in S\Rightarrow a*b\in S$.
  \item Associativity: $a*(b*c)=(a*b)*c$.
  \item Identity: $a*e=e*a=a$.
  \item Inverse: $a*a^{-1}=a^{-1}*a=e$.
\end{itemize}
\end{frame}

% 4
\begin{frame}{A field supplies the scalars of a vector space}
\begin{columns}[T,onlytextwidth]
\column{0.45\textwidth}
\begin{block}{Ring}
A set with addition and multiplication, where addition forms an Abelian group and multiplication distributes over addition.
\end{block}
\begin{exampleblock}{Example}
$(\mathbb Z,+,\cdot)$ is a ring, but not a field.
\end{exampleblock}
\column{0.49\textwidth}
\begin{block}{Field}
A ring in which each nonzero element has a multiplicative inverse.
\end{block}
\begin{exampleblock}{Examples}
$\mathbb Q$, $\mathbb R$ and $\mathbb C$ are fields. Vector spaces in machine learning normally use $\mathbb R$.
\end{exampleblock}
\end{columns}
\end{frame}

% 5
\begin{frame}{Vectors combine magnitude, direction and features}
\begin{columns}[T,onlytextwidth]
\column{0.48\textwidth}
\begin{block}{Vector-space notation}
\[
(V,F,+,\cdot)
\]
\begin{itemize}
  \item $V$: set of vectors
  \item $F$: scalar field
  \item $+$: vector addition
  \item $\cdot$: scalar multiplication
\end{itemize}
\end{block}
\column{0.46\textwidth}
\begin{exampleblock}{Example in $\mathbb R^2$}
\[
\bm v=\begin{bmatrix}2\\3\end{bmatrix},\qquad
3\bm v=\begin{bmatrix}6\\9\end{bmatrix}
\]
The vector can represent a geometric displacement or two numerical features of an object.
\end{exampleblock}
\end{columns}
\end{frame}

% 6
\begin{frame}{Vector operations act component by component}
\begin{columns}[T,onlytextwidth]
\column{0.48\textwidth}
\begin{block}{Addition}
\[
\bm u+\bm v=
\begin{bmatrix}u_1+v_1\\u_2+v_2\\\vdots\\u_n+v_n\end{bmatrix}
\]
\end{block}
\column{0.46\textwidth}
\begin{exampleblock}{Worked example}
\[
\begin{bmatrix}3\\5\\2\end{bmatrix}+
\begin{bmatrix}1\\3\\4\end{bmatrix}=
\begin{bmatrix}4\\8\\6\end{bmatrix}
\]
\end{exampleblock}
\end{columns}
\vspace{0.25cm}
Scalar multiplication stretches, shrinks, or reverses a vector while preserving its line of direction.
\end{frame}

% 7
\begin{frame}{The dot product links algebra and geometry}
\keyeq{\langle\bm u,\bm v\rangle=\bm u^{\mathsf T}\bm v=\sum_{i=1}^{n}u_iv_i}
\begin{columns}[T,onlytextwidth]
\column{0.47\textwidth}
\begin{block}{Magnitude}
\[
\|\bm v\|=\sqrt{\langle\bm v,\bm v\rangle}
\]
The squared norm is $\|\bm v\|^2=\langle\bm v,\bm v\rangle$.
\end{block}
\column{0.47\textwidth}
\begin{exampleblock}{Example}
For $\bm v=(3,4)$,
\[
\|\bm v\|=\sqrt{3^2+4^2}=5.
\]
The corresponding unit vector is $(3/5,4/5)$.
\end{exampleblock}
\end{columns}
\end{frame}

% 8
\begin{frame}{Distance and angle follow from the inner product}
\begin{columns}[T,onlytextwidth]
\column{0.46\textwidth}
\begin{block}{Euclidean distance}
\[
d(\bm u,\bm v)=\|\bm v-\bm u\|
\]
For $\bm u=(1,2)$ and $\bm v=(4,6)$,
\[
d=\sqrt{3^2+4^2}=5.
\]
\end{block}
\column{0.48\textwidth}
\begin{exampleblock}{Angle between nonzero vectors}
\[
\cos\theta=\frac{\bm u\cdot\bm v}{\|\bm u\|\,\|\bm v\|}
\]
\[
\theta=\cos^{-1}\!\left(\frac{\bm u\cdot\bm v}{\|\bm u\|\,\|\bm v\|}\right)
\]
\end{exampleblock}
\end{columns}
\end{frame}

% 9
\begin{frame}{A metric formalizes what distance must mean}
A metric space is a pair $(X,d)$, where $X$ is a set and $d$ assigns a distance to every pair of points.
\begin{columns}[T,onlytextwidth]
\column{0.47\textwidth}
\begin{block}{Essential properties}
\begin{enumerate}
  \item $d(\bm x,\bm y)\ge 0$
  \item $d(\bm x,\bm y)=0\iff\bm x=\bm y$
  \item $d(\bm x,\bm y)=d(\bm y,\bm x)$
\end{enumerate}
\end{block}
\column{0.47\textwidth}
\begin{exampleblock}{Triangle inequality}
\[
d(\bm x,\bm z)\le d(\bm x,\bm y)+d(\bm y,\bm z)
\]
In Euclidean space, the usual metric is $d(\bm u,\bm v)=\|\bm v-\bm u\|$.
\end{exampleblock}
\end{columns}
\end{frame}

% 10
\begin{frame}{Orthogonality is detected by a zero dot product}
\begin{columns}[T,onlytextwidth]
\column{0.47\textwidth}
\begin{block}{Perpendicular vectors}
\[
\bm u\perp\bm v\quad\Longleftrightarrow\quad \bm u\cdot\bm v=0
\]
Example:
\[
(1,2)\cdot(2,-1)=2-2=0.
\]
\end{block}
\column{0.47\textwidth}
\begin{exampleblock}{Cosine similarity}
\[
\operatorname{sim}(\bm u,\bm v)=
\frac{\bm u\cdot\bm v}{\|\bm u\|\,\|\bm v\|}
\]
$1$: same direction\\
$0$: orthogonal\\
$-1$: opposite direction
\end{exampleblock}
\end{columns}
\end{frame}

% 11
\begin{frame}{A hyperplane generalizes a line and a plane}
\keyeq{\bm w^{\mathsf T}\bm x+b=0}
\begin{columns}[T,onlytextwidth]
\column{0.46\textwidth}
\begin{block}{By dimension}
\begin{itemize}
  \item In $\mathbb R^2$: a line
  \item In $\mathbb R^3$: a plane
  \item In $\mathbb R^n$: an $(n-1)$-dimensional hyperplane
\end{itemize}
\end{block}
\column{0.48\textwidth}
\begin{exampleblock}{Parameter roles}
\begin{itemize}
  \item $\bm w$: normal or weight vector
  \item $b$: bias or intercept
  \item Changing $\bm w$ changes orientation
  \item Changing $b$ shifts the boundary
\end{itemize}
\end{exampleblock}
\end{columns}
\end{frame}

% 12
\begin{frame}{Why $(a,b)$ is perpendicular to $ax+by+c=0$}
Choose any two points $P=(x_1,y_1)$ and $Q=(x_2,y_2)$ on the line. Then
\[
ax_1+by_1+c=0,
\qquad
ax_2+by_2+c=0.
\]
Subtracting the second equation from the first gives
\[
a(x_1-x_2)+b(y_1-y_2)=0.
\]
But the direction vector along the line is $\overrightarrow{QP}=(x_1-x_2,y_1-y_2)$, so
\[
(a,b)\cdot\overrightarrow{QP}=0.
\]
\begin{exampleblock}{Conclusion}
A zero dot product means perpendicular vectors. Therefore, the coefficient vector $(a,b)$ is normal to the line $ax+by+c=0$.
\end{exampleblock}
\end{frame}

% 13
\begin{frame}{The weight vector is normal to the hyperplane}
Let $\bm x_1$ and $\bm x_2$ be two points on the same hyperplane:
\[
\bm w^{\mathsf T}\bm x_1+b=0,
\qquad
\bm w^{\mathsf T}\bm x_2+b=0.
\]
Subtracting the equations gives:
\[
\bm w^{\mathsf T}(\bm x_1-\bm x_2)=0.
\]
\begin{exampleblock}{Geometric meaning}
$\bm x_1-\bm x_2$ lies along the hyperplane. Its dot product with $\bm w$ is zero, so $\bm w$ is perpendicular to every direction along the hyperplane.
\end{exampleblock}
\end{frame}

% 14
\begin{frame}{Augmented vectors absorb the bias term}
\begin{columns}[T,onlytextwidth]
\column{0.45\textwidth}
\[
\widetilde{\bm w}=
\begin{bmatrix}\bm w\\b\end{bmatrix},
\qquad
\widetilde{\bm x}=
\begin{bmatrix}\bm x\\1\end{bmatrix}
\]
\column{0.49\textwidth}
\[
\widetilde{\bm w}^{\mathsf T}\widetilde{\bm x}
=\bm w^{\mathsf T}\bm x+b
\]
\end{columns}
\vspace{0.45cm}
\begin{block}{Two useful forms}
The homogeneous form $\bm w^{\mathsf T}\bm x=0$ passes through the origin. The affine form $\bm w^{\mathsf T}\bm x+b=0$ can be shifted away from the origin.
\end{block}
\end{frame}

% 15
\begin{frame}{A perceptron classifies by choosing a side}
\keyeq{z=\bm w^{\mathsf T}\bm x+b\qquad\hat y=\phi(z)}
\begin{columns}[T,onlytextwidth]
\column{0.47\textwidth}
\begin{block}{Step activation}
\[
\phi(z)=
\begin{cases}
1,&z\ge 0,\\
0,&z<0.
\end{cases}
\]
\end{block}
\column{0.47\textwidth}
\begin{exampleblock}{Decision boundary}
\[
\bm w^{\mathsf T}\bm x+b=0
\]
Positive scores belong to one class; negative scores belong to the other.
\end{exampleblock}
\end{columns}
\end{frame}

% 16
\begin{frame}{A complete classification takes one dot product}
Consider the decision boundary:
\[
2x_1+x_2-4=0,
\qquad \bm w=(2,1),\quad b=-4.
\]
\begin{columns}[T,onlytextwidth]
\column{0.46\textwidth}
\begin{block}{Input $\bm x=(3,1)$}
\[
z=2(3)+1-4=3
\]
Since $z>0$, the prediction is $\hat y=1$.
\end{block}
\column{0.46\textwidth}
\begin{exampleblock}{Input $\bm x=(1,1)$}
\[
z=2(1)+1-4=-1
\]
Since $z<0$, the prediction is $\hat y=0$.
\end{exampleblock}
\end{columns}
\end{frame}

% 17
\begin{frame}{Learning moves the boundary after an error}
For labels $y\in\{0,1\}$, a common update is:
\[
\bm w\leftarrow\bm w+\eta(y-\hat y)\bm x,
\qquad
b\leftarrow b+\eta(y-\hat y).
\]
\begin{itemize}
  \item $\eta$ is the learning rate.
  \item A correct prediction gives $y-\hat y=0$, so no update is made.
  \item A wrong prediction moves the separating hyperplane toward a better position.
\end{itemize}
\begin{alertblock}{Limitation}
A single perceptron learns only linearly separable patterns. It can represent AND and OR, but not XOR.
\end{alertblock}
\end{frame}

% 18
\begin{frame}{The mathematics converges in the perceptron equation}
\begin{itemize}
  \item Algebraic structures provide the rules behind vector spaces.
  \item Vectors encode direction, magnitude and numerical features.
  \item Dot products produce norms, angles and orthogonality tests.
  \item A weight vector determines a hyperplane's orientation.
  \item Bias shifts the hyperplane, while an activation function assigns the class.
\end{itemize}
\vspace{0.35cm}
\keyeq{\boxed{\hat y=\phi(\bm w^{\mathsf T}\bm x+b)}}
\begin{center}
\color{muted}This equation connects linear algebra directly to binary classification.
\end{center}
\end{frame}


% 20
\begin{frame}[plain]
\begin{center}
\vspace{1.3cm}
  {\color{navy}\rule{2.0cm}{1.5pt}}\\[0.55cm]
{\fontsize{34}{38}\selectfont\bfseries\color{navy}Thank You}\\[0.45cm]
\end{center}
\end{frame}

\end{document}
